\chapter[Localization and Toric Riemann--Roch]
{Localization and Toric\\ Riemann--Roch}
\label{chap:localization-toric-rr}

\begin{chapteroverview}
Chapter~8 showed that the fan of a smooth complete toric variety
determines its cohomology ring. We now ask what we can \emph{do} with
that ring. The answer will be unexpectedly arithmetic: characteristic
classes in cohomology compute Euler characteristics of line bundles,
and on a projective toric variety those Euler characteristics become
counts of lattice points.

There is a second idea running through the chapter. Ordinary
cohomology forgets the torus action, even though that action is the
source of most of the geometry. Equivariant cohomology remembers it.
For a smooth toric variety it has a strikingly combinatorial model as
piecewise polynomial functions on the fan, and localization reduces
global integrals to contributions from torus fixed points.

The chapter follows the progression
\[
\boxed{
\begin{gathered}
\text{Chern character and Todd class}
\longrightarrow
\text{piecewise polynomial equivariant classes}
\\[-1mm]
\longrightarrow
\text{fixed-point localization}
\longrightarrow
\text{toric Riemann--Roch and lattice points}.
\end{gathered}}
\]
Throughout, varieties are defined over \(\CC\). Unless stated
otherwise, the toric varieties in this chapter are smooth and complete.
\end{chapteroverview}

\begin{threedictionaries}
The same toric data will now appear in three new roles.

\medskip

\begin{center}
\small
\begin{tabularx}{\linewidth}{@{}
  >{\raggedright\arraybackslash}X
  >{\raggedright\arraybackslash}X
  >{\raggedright\arraybackslash}X@{}}
\toprule
\textbf{Fan and polytope}
&
\textbf{Equivariant topology}
&
\textbf{Index theory}
\\
\midrule
ray \(\rho\) and divisor \(D_\rho\)
&
degree-two class
&
factor in \(\operatorname{Td}(X)\)
\\
polynomial on each cone
&
equivariant cohomology class
&
fixed-point restriction
\\
maximal cone \(\sigma\)
&
torus fixed point \(x_\sigma\)
&
localization contribution
\\
lattice points of \(P_D\)
&
weights of global sections
&
Euler characteristic of \(\mathcal O(D)\)
\\
\bottomrule
\end{tabularx}
\end{center}
\end{threedictionaries}

\begin{surfacelaboratory}
On a smooth complete toric surface, the theory developed in this
chapter will collapse to an especially concrete formula. For a divisor
\(D\),
\[
\boxed{
\chi\bigl(X,\mathcal O_X(D)\bigr)
=
1+\frac12 D\cdot(D-K_X).
}
\]
Thus the intersection form recovered from the fan in the challenge at
the end of Chapter~8 already contains enough information to compute
Euler characteristics of every line bundle. The first section will
derive this formula; the rest of the chapter will explain why the same
phenomenon in every dimension can be localized at the fixed points and,
for ample divisors, read as lattice-point counting.
\end{surfacelaboratory}

\section{Chern characters, Todd classes, and Riemann--Roch}
\label{sec:chern-todd-hrr}
\index{Chern class}
\index{Chern character}
\index{Todd class}
\index{Hirzebruch--Riemann--Roch}

The cohomology ring from Chapter~8 records intersections of divisor
classes. Riemann--Roch asks us to package those classes in a slightly
different way. Given a line bundle \(L\), we want to recover the
alternating sum
\[
\chi(X,L)
:=
\sum_{i=0}^{\dim X}(-1)^i\dim_\CC H^i(X,L)
\]
from cohomological data on \(X\).

Only three ingredients will be needed: Chern classes, the Chern
character, and the Todd class. We introduce only the parts of their
general theory that enter the toric calculation.

\subsection{The tangent bundle from the boundary divisors}

Let \(E\) be a complex vector bundle of rank \(r\). Its Chern classes
are classes
\[
c_i(E)\in H^{2i}(X;\ZZ),
\qquad
c(E)=1+c_1(E)+\cdots+c_r(E).
\]
For us, the two essential rules are
\[
c(L)=1+c_1(L)
\]
for a line bundle \(L\), and the Whitney product formula
\[
c(E)=c(E')c(E'')
\]
for a short exact sequence
\[
0\longrightarrow E'
\longrightarrow E
\longrightarrow E''
\longrightarrow0.
\]
If \(D\) is a Cartier divisor, then
\[
c_1\bigl(\mathcal O_X(D)\bigr)=[D].
\]

For a smooth toric variety, there is a generalized Euler sequence that
plays the role of the familiar Euler sequence on projective space. If
\(r=|\Sigma(1)|\) and \(n=\dim X_\Sigma\), its tangent-bundle form is
\[
0
\longrightarrow
\mathcal O_{X_\Sigma}^{\,r-n}
\longrightarrow
\bigoplus_{\rho\in\Sigma(1)}
\mathcal O_{X_\Sigma}(D_\rho)
\longrightarrow
T_{X_\Sigma}
\longrightarrow0.
\]
It is the dual of Cox--Little--Schenck, Theorem~8.1.6. We need only
this consequence of the sequence here.

The trivial bundle on the left has total Chern class \(1\), so the
Whitney formula immediately makes the Chern classes combinatorial.

\begin{proposition}[Chern class of a smooth toric variety]
\label{prop:toric-total-chern-class}
Let \(X_\Sigma\) be smooth and complete. Then
\[
\boxed{
c(T_{X_\Sigma})
=
\prod_{\rho\in\Sigma(1)}\bigl(1+[D_\rho]\bigr).
}
\]
In particular,
\[
c_1(T_{X_\Sigma})
=
\sum_\rho[D_\rho]
=
[-K_{X_\Sigma}],
\]
and, if \(n=\dim X_\Sigma\),
\[
c_n(T_{X_\Sigma})
=
|\Sigma(n)|[\mathrm{pt}].
\]
\end{proposition}

\noindent
\textit{Reference.}
This is Cox--Little--Schenck, Proposition~13.1.2. Their
Example~13.1.1 derives the same formula from the cotangent form of the
generalized Euler sequence.

\begin{proof}[Conceptual proof sketch]
Apply the Whitney product formula to the tangent Euler sequence. Each
summand \(\mathcal O(D_\rho)\) contributes \(1+[D_\rho]\), while the
trivial bundle contributes nothing. The degree-two part of the product
is \(\sum_\rho[D_\rho]=[-K_X]\).

For the top class, expand the product in the cohomology ring from
Chapter~8. A product of \(n\) distinct invariant-divisor classes is the
class of a point precisely when the corresponding rays form a maximal
cone, and is zero otherwise. Thus every maximal cone contributes one
copy of \([\mathrm{pt}]\).
\end{proof}

The last statement reconnects immediately with Chapter~8:
\[
\int_X c_n(T_X)
=
|\Sigma(n)|
=
\chi_{\mathrm{top}}(X).
\]
The topological Euler characteristic counted fixed points. The
holomorphic Euler characteristics we want next require a different
combination of the same Chern data.

\subsection{The Chern character is an exponential}

For a line bundle there is almost nothing to remember.

\begin{definition}[Chern character of a line bundle]
\label{def:chern-character-line-bundle}
Let \(L\) be a line bundle on a smooth complete variety. Its
\emph{Chern character} is
\[
\ch(L)
:=
e^{c_1(L)}
=
1+c_1(L)+\frac{c_1(L)^2}{2!}
+\frac{c_1(L)^3}{3!}+\cdots
\in H^*(X;\QQ).
\]
\end{definition}

The apparent infinite series is actually finite on an
\(n\)-dimensional variety: every term above cohomological degree
\(2n\) vanishes. Thus, for \(L=\mathcal O_X(D)\),
\[
\boxed{
\ch\bigl(\mathcal O_X(D)\bigr)=e^{[D]}.
}
\]
The exponential is exactly the right operation because tensor product
of line bundles becomes multiplication:
\[
\ch(L_1\otimes L_2)
=
\ch(L_1)\ch(L_2).
\]

\begin{remark}
The denominators are the first reason rational cohomology naturally
appears in Riemann--Roch. Nothing is being completed analytically:
all power series in this chapter are truncated by dimension.
\end{remark}

\subsection{The Todd class is the correction factor}

Suppose formally that the Chern classes of \(T_X\) factor as
\[
c(T_X)=\prod_{i=1}^n(1+\xi_i).
\]
The symbols \(\xi_i\) are called the \emph{Chern roots}. The splitting
principle allows us to calculate with them as though the tangent bundle
were a sum of line bundles. Define
\[
Q(x)
:=
\frac{x}{1-e^{-x}}
=
1+\frac{x}{2}+\frac{x^2}{12}
-\frac{x^4}{720}+\cdots.
\]

\begin{definition}[Todd class]
\label{def:todd-class}
For a smooth complete \(n\)-fold \(X\), the \emph{Todd class} is
\[
\operatorname{Td}(X)
:=
\prod_{i=1}^n Q(\xi_i)
=
\prod_{i=1}^n
\frac{\xi_i}{1-e^{-\xi_i}}
\in H^*(X;\QQ).
\]
\end{definition}

The first terms, written without Chern roots, are
\[
\operatorname{Td}(X)
=
1+\frac12c_1
+\frac1{12}(c_1^2+c_2)
+\frac1{24}c_1c_2+\cdots.
\]
As always, only terms through degree \(2\dim X\) matter.

For a toric variety the formal Chern roots can disappear entirely.

\begin{theorem}[The toric Todd class]
\label{thm:toric-todd-class}
Let \(X_\Sigma\) be a smooth complete toric variety. Then
\[
\boxed{
\operatorname{Td}(X_\Sigma)
=
\prod_{\rho\in\Sigma(1)}
\frac{[D_\rho]}{1-e^{-[D_\rho]}}.
}
\]
The right-hand side is interpreted through the power series
\(Q(x)\), so every factor has constant term \(1\).
\end{theorem}

\noindent
\textit{Reference.}
This is Cox--Little--Schenck, Theorem~13.1.6.

\begin{proof}[Conceptual proof sketch]
Characteristic classes defined by a power series are multiplicative
for exact sequences. Apply this to the tangent Euler sequence. The
trivial summands have Todd class \(1\), while
\(\mathcal O(D_\rho)\) contributes
\[
Q([D_\rho])
=
\frac{[D_\rho]}{1-e^{-[D_\rho]}}.
\]
Multiplying over the rays gives the formula. Equivalently,
\cref{prop:toric-total-chern-class} says that the divisor classes
behave, for symmetric characteristic-class calculations, like a
redundant list of Chern roots.
\end{proof}

This is the first major payoff of the toric boundary. A characteristic
class of the tangent bundle has become a product indexed by the rays of
the fan, and every factor lives in the cohomology ring we already know
how to compute.

\subsection{Hirzebruch--Riemann--Roch}

We can now state the theorem that motivates the chapter. If
\(\alpha\in H^*(X;\QQ)\) and \(\dim X=n\), write
\[
\int_X\alpha
\]
for evaluation of the degree-\(2n\) part of \(\alpha\) on the
fundamental class of \(X\). In particular,
\[
\int_X[\mathrm{pt}]=1.
\]

\begin{theorem}[Hirzebruch--Riemann--Roch for toric line bundles]
\label{thm:toric-hrr}
Let \(X=X_\Sigma\) be a smooth complete toric variety and let \(L\) be
a line bundle on \(X\). Then
\[
\boxed{
\chi(X,L)
=
\int_X\ch(L)\operatorname{Td}(X).
}
\]
For \(L=\mathcal O_X(D)\), this becomes
\[
\boxed{
\chi\bigl(X,\mathcal O_X(D)\bigr)
=
\int_X
e^{[D]}
\prod_{\rho\in\Sigma(1)}
\frac{[D_\rho]}{1-e^{-[D_\rho]}}.
}
\]
\end{theorem}

\noindent
\textit{Reference.}
The general HRR formula is equation~(13.0.1) in
Cox--Little--Schenck. The line-bundle statement for smooth complete
toric varieties is their Theorem~13.3.8, obtained from equivariant
Riemann--Roch by taking the nonequivariant limit.

\begin{proof}[Conceptual proof sketch]
We will not reproduce the general proof of Riemann--Roch. Its message
is that the Chern character converts the algebra of a line bundle into
cohomology, while the Todd class corrects for the geometry of the
tangent bundle. Their product has a top-degree component, and
integration of that component recovers the alternating sum of sheaf
cohomology dimensions.

For toric varieties there is a more concrete mechanism behind this
global statement. In Sections~9.2 and~9.3 we will retain the torus
action in equivariant cohomology and then use localization to reduce an
integral to fixed-point contributions. Taking the nonequivariant limit
of that calculation is precisely the route used in CLS to prove the
toric statement above.
\end{proof}

\paragraph{A four-step computation.}
For a smooth complete toric variety, HRR has become a calculation in
the ring from Section~8.4:
\begin{enumerate}[label=\arabic*.,wide=0pt]
\item write \([D]\) and the invariant classes \([D_\rho]\) in
      \(H^*(X;\QQ)\);
\item expand \(e^{[D]}\) only through degree \(2\dim X\);
\item expand each Todd factor
      \(1+\frac12[D_\rho]+\frac1{12}[D_\rho]^2+\cdots\);
\item use the fan relations and keep only the top-degree coefficient.
\end{enumerate}
That coefficient is \(\chi(X,\mathcal O_X(D))\).

\subsection{The first calculation: line bundles on
\texorpdfstring{\(\PP^1\)}{P1}}

Let \(H\in H^2(\PP^1;\ZZ)\) be the class of a point divisor. From
Chapter~8,
\[
H^*(\PP^1;\ZZ)
\simeq
\frac{\ZZ[H]}{\langle H^2\rangle},
\qquad
\int_{\PP^1}H=1.
\]
There are two invariant divisors, both of class \(H\). Since
\(H^2=0\), the Todd product truncates immediately:
\[
\operatorname{Td}(\PP^1)
=
\left(1+\frac12H\right)^2
=
1+H.
\]
For \(L=\mathcal O_{\PP^1}(d)\),
\[
\ch(L)=e^{dH}=1+dH.
\]
Therefore
\[
\begin{aligned}
\chi\bigl(\PP^1,\mathcal O(d)\bigr)
&=
\int_{\PP^1}(1+dH)(1+H)\\
&=
\int_{\PP^1}(d+1)H
=d+1.
\end{aligned}
\]

Notice that HRR computes the \emph{Euler characteristic}, not merely
the number of global sections. When \(d\ge0\), higher cohomology
vanishes and this says
\[
h^0\bigl(\PP^1,\mathcal O(d)\bigr)=d+1.
\]
For \(d=-2\), however,
\(\chi(\mathcal O(-2))=-1\): the negative sign records a nonzero
\(H^1\), not a negative-dimensional space of sections.

\subsection{Surfaces: Riemann--Roch becomes intersection theory}

Let \(X\) be a smooth complete surface. Since \(c_1(T_X)=-K_X\),
\[
\operatorname{Td}(X)
=
1-\frac12K_X
+\frac1{12}(K_X^2+c_2(X)),
\]
while
\[
\ch\bigl(\mathcal O_X(D)\bigr)
=
1+D+\frac12D^2.
\]
Only cohomological degree four survives integration. HRR therefore
gives
\[
\chi\bigl(X,\mathcal O_X(D)\bigr)
=
\frac12D\cdot(D-K_X)
+\frac1{12}\int_X(K_X^2+c_2(X)).
\]
Applying the same formula to \(D=0\) identifies the last term with
\(\chi(X,\mathcal O_X)\). Hence
\[
\boxed{
\chi\bigl(X,\mathcal O_X(D)\bigr)
=
\chi(X,\mathcal O_X)
+\frac12D\cdot(D-K_X).
}
\]
This is the familiar surface Riemann--Roch formula. For a smooth
complete toric variety,
\[
H^i(X,\mathcal O_X)=0\qquad(i>0),
\]
so \(\chi(X,\mathcal O_X)=1\). We recover the formula announced at
the beginning of the chapter:
\[
\boxed{
\chi\bigl(X,\mathcal O_X(D)\bigr)
=
1+\frac12D\cdot(D-K_X).
}
\]

\noindent
\textit{Reference.}
The surface calculation is Cox--Little--Schenck,
Example~13.1.7 and Exercise~13.1.2; see their equation~(13.1.7).

\subsection{A triangle already knows Riemann--Roch}

Take \(X=\PP^2\) and let \(H\) be the hyperplane class. We know
\[
H^*(\PP^2;\ZZ)
\simeq
\frac{\ZZ[H]}{\langle H^3\rangle},
\qquad
\int_{\PP^2}H^2=1.
\]
The three invariant divisors all represent \(H\), so
\[
\begin{aligned}
\operatorname{Td}(\PP^2)
&=
\left(1+\frac12H+\frac1{12}H^2\right)^3\\
&=
1+\frac32H+H^2.
\end{aligned}
\]
For \(L=\mathcal O_{\PP^2}(d)\),
\[
\ch(L)
=
1+dH+\frac{d^2}{2}H^2.
\]
Multiplying and taking the coefficient of \(H^2\),
\[
\begin{aligned}
\chi\bigl(\PP^2,\mathcal O(d)\bigr)
&=
1+\frac32d+\frac12d^2\\
&=
\frac{(d+1)(d+2)}{2}.
\end{aligned}
\]
For \(d\ge0\), higher cohomology vanishes, and the last number is
also
\[
\#\bigl(d\Delta_2\cap\ZZ^2\bigr),
\]
the number of lattice points in the degree-\(d\) triangle. At the
moment this is a familiar coincidence: monomials of degree \(d\)
already explain the same count. By the end of the chapter, it will be
the model for a general toric Riemann--Roch principle.

\begin{dictionarytable}
Section~8.4 turned fan data into a cohomology ring. The present section
has added two universal operations inside that ring:
\[
[D]
\longmapsto
e^{[D]},
\qquad
\{[D_\rho]\}_\rho
\longmapsto
\prod_\rho
\frac{[D_\rho]}{1-e^{-[D_\rho]}}.
\]
Their top-degree product is an integer: the Euler characteristic of a
line bundle. What is still hidden is the torus action itself. The next
section keeps that information and turns cohomology classes into
polynomials living directly on the fan.
\end{dictionarytable}

\begin{guidingquestion}
Ordinary cohomology has forgotten the torus action. Is there a version
of cohomology in which a class can be read cone by cone from the fan,
and whose restrictions to maximal cones remember the torus fixed
points?
\end{guidingquestion}

\section{Equivariant cohomology and piecewise polynomials}
\label{sec:equivariant-piecewise-polynomials}
\index{equivariant cohomology}
\index{Borel construction}
\index{piecewise polynomial}
\index{support function!as an equivariant class}

There is a small paradox in the cohomology ring from Chapter~8. The
torus action organizes nearly everything we do with a toric variety,
yet the ordinary cohomology ring does not remember the action itself.
Equivariant cohomology repairs this by keeping track of a space
\emph{together with} its torus action.

For the rest of this section we write
\[
K_N=\Hom_{\ZZ}(M,S^1)\simeq(S^1)^n
\]
for the compact torus from Section~8.1. The algebraic torus
\(T_N\simeq(\CC^*)^n\) deformation retracts onto \(K_N\), so this is the
same topological equivariant theory that Cox--Little--Schenck denote by
\(H_T^*(-)\). We use rational coefficients in the main statements;
for a smooth fan the same combinatorics also has an integral version.

\subsection{The Borel construction: remembering the action}

If a compact group \(K\) acts freely on a space \(X\), the quotient
\(X/K\) is often the right space to study. A torus action on a toric
variety is very far from free: fixed points are central to the story.
The Borel construction replaces the action by an auxiliary free one.

Choose a contractible space \(EK\) on which \(K\) acts freely, and put
\[
BK:=EK/K.
\]
The diagonal action on \(EK\times X\) is free, so we may form
\[
EK\times_KX:=(EK\times X)/K.
\]

\begin{definition}[Equivariant cohomology]
\label{def:equivariant-cohomology}
The \emph{\(K\)-equivariant cohomology} of a \(K\)-space \(X\) is
\[
H_K^*(X;\QQ)
:=
H^*(EK\times_KX;\QQ).
\]
\end{definition}

\begin{agcheatsheet}
For an algebraic geometer, the easiest way to remember equivariant
cohomology is as a cohomology theory for a space \emph{together with an
action}.  Write
\[
X_K:=EK\times_KX
\]
for the Borel construction, sometimes called the \emph{homotopy
quotient}.  It is designed to satisfy the following basic tests.
An equivariant map \(f:X\to Y\) induces a map \(X_K\to Y_K\), so
\(H_K^*(-)\) is a contravariant functor on \(K\)-spaces, just as
ordinary cohomology is a contravariant functor on spaces.

\begin{enumerate}[label=\arabic*.,wide=0pt]
\item \textbf{Free actions give ordinary quotients.}
If \(K\) acts freely on \(X\), then
\[
X_K\simeq X/K,
\qquad
H_K^*(X)\simeq H^*(X/K).
\]
Thus equivariant cohomology really does extend the familiar operation
of taking the cohomology of a geometric quotient.

\item \textbf{A trivial action gives a tensor product.}
If \(K\) acts trivially on \(X\), then
\[
X_K=BK\times X,
\]
so, over \(\QQ\), the K\"unneth theorem gives
\[
H_K^*(X;\QQ)
\simeq
H^*(X;\QQ)\otimes_\QQ H_K^*(\mathrm{pt};\QQ).
\]
The second factor is the extra cohomology contributed by the group.

\item \textbf{Equivariant cohomology is a module over the cohomology
of a point.}
The projection \(X_K\to BK\) makes
\[
H_K^*(X;\QQ)
\]
an algebra, and in particular a module, over
\[
H_K^*(\mathrm{pt};\QQ)=H^*(BK;\QQ).
\]
For a torus this coefficient ring is a polynomial ring.  Its variables
are the \emph{equivariant parameters} that keep track of the action.

\item \textbf{There is a nonequivariant limit.}
The fibration
\[
X\longrightarrow X_K\longrightarrow BK
\]
gives a forgetful map \(H_K^*(X)\to H^*(X)\).  Algebraically, forgetting
the action means sending every positive-degree class in \(H^*(BK)\) to
zero.  When \(X\) is \emph{equivariantly formal}---in particular, when
\(H^{\mathrm{odd}}(X;\QQ)=0\)---this gives
\[
H^*(X;\QQ)
\simeq
H_K^*(X;\QQ)
\otimes_{H^*(BK;\QQ)}\QQ,
\]
where \(H^*(BK;\QQ)\to\QQ\) kills all positive-degree elements.

Smooth complete toric varieties are equivariantly formal.  In that
case, as a module over \(H^*(BK;\QQ)\), one may think additively of
\[
H_K^*(X;\QQ)
\simeq
H^*(X;\QQ)\otimes_\QQ H^*(BK;\QQ).
\]
This module isomorphism is not canonical in general, and the tensor
product formula should \emph{not} be taken as a property of every
group action.
\end{enumerate}

These four tests are usually enough intuition for the calculations in
this chapter: equivariant cohomology interpolates between the
cohomology of a quotient and the ordinary cohomology of \(X\), while
remembering the group action through the coefficient ring \(H^*(BK)\).
\end{agcheatsheet}

\noindent
\textit{Reference.}
This is Cox--Little--Schenck, Definition~12.4.6. Their
equations~(12.4.7)--(12.4.9) compute the coefficient ring for a torus.

\begin{example}[The coefficient ring for \(S^1\)]
\label{ex:bs1-equivariant-point}
Let \(S^1\) act on the infinite-dimensional unit sphere
\(S^\infty\subset\CC^\infty\) by scalar multiplication.  This action
is free, and \(S^\infty\) is contractible, so one may take
\[
ES^1=S^\infty,
\qquad
BS^1=\PP^\infty,
\]
where \(\PP^\infty\) is the direct limit of the finite-dimensional
projective spaces.  Since
\[
H^*(BS^1;\ZZ)\simeq\ZZ[t],
\qquad
t=c_1\bigl(\mathcal O_{\PP^\infty}(1)\bigr),
\qquad
\deg t=2,
\]
we obtain
\[
\boxed{H_{S^1}^*(\mathrm{pt};\ZZ)\simeq\ZZ[t].}
\]
The finite-dimensional approximations are already familiar to an
algebraic geometer:
\[
S^{2r+1}/S^1=\PP^r.
\]
As \(r\) grows, their cohomology stabilizes in every fixed degree and
produces the polynomial ring above.
\end{example}

\begin{example}[The algebraic torus]
\label{ex:algebraic-torus-equivariant-point}
Now let
\[
T_N=\Hom_\ZZ(M,\CC^*)\simeq(\CC^*)^n
\]
be the algebraic torus.  Over \(\CC\), it deformation retracts onto its
maximal compact torus
\[
K_N=\Hom_\ZZ(M,S^1)\simeq(S^1)^n.
\]
Consequently their classifying spaces have the same homotopy type, and
taking products of the preceding example gives
\[
BT_N\simeq BK_N\simeq(\PP^\infty)^n.
\]

More intrinsically, every character \(m\in M\) gives a homomorphism
\(T_N\to\CC^*\), hence a map of classifying spaces to
\(B\CC^*\simeq\PP^\infty\).  Pulling back the universal degree-two
class identifies the character lattice with the degree-two part of the
coefficient ring.  Thus
\[
\boxed{
H_{T_N}^*(\mathrm{pt};\ZZ)
\simeq
H_{K_N}^*(\mathrm{pt};\ZZ)
\simeq
\operatorname{Sym}_{\ZZ}(M),
}
\]
where elements of \(M\) have cohomological degree \(2\).  After
tensoring with \(\QQ\), we will write
\[
S:=\operatorname{Sym}_{\QQ}(M_\QQ).
\]
If a basis \(m_1,\ldots,m_n\) of \(M\) is chosen, this simply reads
\[
S\simeq\QQ[t_1,\ldots,t_n],
\qquad \deg t_i=2.
\]
The same approximation can be written entirely in algebraic geometry.
After choosing coordinates on \(T_N\), let it act factor by factor on
\[
U_r:=\bigl(\AAA^{r+1}\setminus\{0\}\bigr)^n.
\]
The action is free and
\[
U_r/T_N\simeq(\PP^r)^n.
\]
As \(r\to\infty\), these quotients approximate \(BT_N\) in increasing
cohomological degrees.  Thus the topological classifying space above
is already visible through a familiar sequence of algebraic quotients.
\end{example}

Thus equivariant cohomology is not merely a graded \(\QQ\)-algebra: it
is naturally an \(S\)-algebra. The extra variables in \(S\) are the
parameters that remember the torus action.

\begin{pauseandthink}
Test the cheatsheet on \(S^1\) itself.  If \(S^1\) acts on itself by
rotation, the action is free, so what is \(H_{S^1}^*(S^1;\QQ)\)?  What
does it become if the action is trivial instead?  The underlying space
has not changed; only the action has.  Where does equivariant
cohomology record the difference?
\end{pauseandthink}

\subsection{What a toric equivariant class looks like}

Now take a smooth complete fan \(\Sigma\subseteq N_\RR\). Every maximal
cone \(\sigma\in\Sigma(n)\) corresponds to a torus fixed point
\(x_\sigma\). Restriction to the fixed points gives a map
\[
H_{K_N}^*(X_\Sigma;\QQ)
\longrightarrow
\bigoplus_{\sigma\in\Sigma(n)}
H_{K_N}^*(x_\sigma;\QQ)
\simeq
\bigoplus_{\sigma\in\Sigma(n)}S.
\]
So an equivariant class determines one polynomial \(f_\sigma\in S\)
at each maximal cone. Those polynomials cannot be chosen
independently: neighboring fixed points are joined by one-dimensional
orbits, and this forces their polynomials to agree on the common face.

This is exactly the gluing rule for piecewise polynomial functions.

\begin{definition}[Piecewise polynomial functions on a fan]
\label{def:piecewise-polynomial-fan}
Let \(PP^*(\Sigma;\QQ)\) be the ring of functions
\[
f:|\Sigma|\longrightarrow\RR
\]
whose restriction to every cone is given by a polynomial with rational
coefficients. Equivalently, for a complete fan it is the ring of collections
\[
\{f_\sigma\in S:\sigma\in\Sigma(n)\}
\]
such that
\[
f_\sigma|_{\sigma\cap\sigma'}
=
f_{\sigma'}|_{\sigma\cap\sigma'}
\]
for every pair of maximal cones. A polynomial of ordinary degree \(k\)
is placed in cohomological degree \(2k\).
\end{definition}

For instance, degree-two equivariant classes correspond to rational
functions that are linear on each cone. The integral classes recover
precisely the integral support functions that first appeared in
Chapter~4.

\begin{theorem}[Equivariant cohomology is piecewise polynomial]
\label{thm:equivariant-piecewise-polynomial}
Let \(X_\Sigma\) be a smooth complete toric variety. Then restriction
to the torus fixed points induces an isomorphism of \(S\)-algebras
\[
\boxed{
H_{K_N}^*(X_\Sigma;\QQ)
\simeq
PP^*(\Sigma;\QQ).
}
\]
The same ring also has the Stanley--Reisner presentation
\[
\boxed{
H_{K_N}^*(X_\Sigma;\QQ)
\simeq
\frac{\QQ[x_\rho:\rho\in\Sigma(1)]}{I_\Sigma},
\qquad
\deg x_\rho=2.
}
\]
Here \(x_\rho\) is the equivariant class of \(D_\rho\).
\end{theorem}

\noindent
\textit{Reference.}
The Stanley--Reisner presentation is Cox--Little--Schenck,
Theorem~12.4.14. The piecewise-polynomial interpretation is the
discussion following their exact sequence~(12.4.18); see especially
Example~12.4.18 for \(\PP^2\). They explain immediately before that
discussion that for a smooth fan the same argument may be carried out
integrally.

\begin{proof}[Conceptual proof sketch]
On the affine chart \(U_\sigma\) of a maximal cone, the torus action
equivariantly retracts onto the fixed point \(x_\sigma\). Thus an
equivariant class restricts there to an element of
\(H_{K_N}^*(x_\sigma)\simeq S\), giving the polynomial \(f_\sigma\).
On an overlap, the two restrictions come from the same global class,
so they agree on the common face.

The nontrivial direction is that these compatibility conditions are
enough. The Chang--Skjelbred exact sequence used in CLS says, in this
toric setting, that equivariant cohomology is determined by the fixed
points and the one-dimensional orbits joining them. Hence every
compatible collection glues to an equivariant class. Describing the
same gluing with one generator for each invariant divisor gives the
Stanley--Reisner presentation: nonfaces give exactly the monomial
relations in \(I_\Sigma\).
\end{proof}

This theorem explains something that was only foreshadowed in
Section~8.4. Ordinary cohomology required \emph{two} ideals,
\(I_\Sigma+J_\Sigma\). Equivariant cohomology requires only
\(I_\Sigma\). The linear relations have not vanished; they have become
the coefficient-ring variables that remember the action.

There is one sign convention worth recording. With our convention
\[
\operatorname{div}(\character{m})
=
\sum_\rho\langle m,u_\rho\rangle D_\rho,
\]
the \(S\)-algebra structure on the Stanley--Reisner presentation sends
\[
m
\longmapsto
-\theta_m
=
-\sum_\rho\langle m,u_\rho\rangle x_\rho.
\]
This is Cox--Little--Schenck's convention and is compatible with the
support-function convention from Chapter~4. The sign will disappear
when we pass to ordinary cohomology.

\subsection{The first picture: \texorpdfstring{\(\PP^1\)}{P1}}

Let \(N=M=\ZZ\), and write \(t\) for the degree-two generator of
\[
S=\QQ[t].
\]
The fan of \(\PP^1\) has two maximal cones
\[
\sigma_+=\RR_{\ge0},
\qquad
\sigma_-=\RR_{\le0},
\]
meeting at the origin. A piecewise polynomial is therefore a pair
\((f_+,f_-)\in S^2\) whose values agree at \(0\). Algebraically,
\[
f_+-f_-\in(t).
\]
Hence
\[
\boxed{
H_{K_N}^*(\PP^1;\QQ)
\simeq
\bigl\{(f_+,f_-)\in\QQ[t]^2:
f_+-f_-\in(t)\bigr\}.
}
\]

\begin{figure}[htbp]
\centering
\begin{tikzpicture}[scale=1.0]
  \draw[->,thick] (-4.2,0) -- (4.2,0);
  \fill (0,0) circle (1.8pt);
  \node[below=5pt] at (0,0) {$0$};
  \node[above=7pt] at (-2.5,0) {$f_-(t)$};
  \node[above=7pt] at (2.5,0) {$f_+(t)$};
  \node[below=8pt] at (-2.5,0) {$\sigma_-$};
  \node[below=8pt] at (2.5,0) {$\sigma_+$};
  \node[align=center,DictionaryBlue] at (0,1.0)
    {$f_+(0)=f_-(0)$\\[-1mm]
     \scriptsize compatibility on the common face};
\end{tikzpicture}
\caption{An equivariant class on \(\PP^1\) is a polynomial on each
ray of the fan, with the two polynomials agreeing where the rays meet.}
\label{fig:p1-piecewise-polynomial}
\end{figure}

This description can be rewritten as a familiar quotient ring. Let
\[
\eta:=(0,t).
\]
Every compatible pair is uniquely \(a(t)(1,1)+b(t)\eta\), and
\[
\eta^2=t\eta.
\]
Therefore
\[
H_{K_N}^*(\PP^1;\QQ)
\simeq
\frac{\QQ[t,\eta]}{\langle\eta(\eta-t)\rangle}.
\]
The relation has not said that \(\eta\) is nilpotent: the parameter
\(t\) records the torus weight. Nilpotence appears only after the torus
parameters are forgotten.

\subsection{Degree two: our old support functions return}

The piecewise-polynomial theorem has a particularly satisfying degree-
two interpretation. Let
\[
D=\sum_\rho a_\rho D_\rho
\]
be an invariant Cartier divisor with Cartier data
\(\{m_\sigma\}\). Recall from Section~4.2 that its support function is
\[
\psi_D|_\sigma=\langle m_\sigma,-\rangle,
\qquad
\psi_D(u_\rho)=-a_\rho.
\]

The divisor also has an equivariant first Chern class
\[
[D]_{K_N}\in H_{K_N}^2(X_\Sigma;\QQ).
\]
At the fixed point \(x_\sigma\), its restriction is precisely the
linear form \(m_\sigma\). Thus under
\cref{thm:equivariant-piecewise-polynomial},
\[
\boxed{
[D]_{K_N}
\longleftrightarrow
\psi_D.
}
\]

\noindent
\textit{Reference.}
The equivariant divisor class and its restriction properties are
Cox--Little--Schenck, Proposition~12.4.13. Combining that proposition
with the Cartier data from their Theorem~4.2.8 gives the identification
above.

This is not a new analogy: it is the same object seen in two subjects.
The linear piece \(m_\sigma\) that gives a local equation for the
divisor is also the weight seen by equivariant cohomology at the fixed
point \(x_\sigma\).

\begin{threedictionaries}
For an invariant Cartier divisor \(D\), the following data are three
descriptions of the same local information.

\medskip

\begin{center}
\small
\begin{tabularx}{\linewidth}{@{}
  >{\raggedright\arraybackslash}X
  >{\raggedright\arraybackslash}X
  >{\raggedright\arraybackslash}X@{}}
\toprule
\textbf{Algebraic geometry}
&
\textbf{Fan}
&
\textbf{Equivariant topology}
\\
\midrule
local equation \(\character{-m_\sigma}\)
&
linear piece \(\psi_D|_\sigma=m_\sigma\)
&
restriction \([D]_{K_N}|_{x_\sigma}=m_\sigma\)
\\
Cartier compatibility
&
pieces agree on faces
&
fixed-point restrictions glue
\\
principal character
&
global linear function
&
class from \(H^*(BK_N)\)
\\
\bottomrule
\end{tabularx}
\end{center}
\end{threedictionaries}

\subsection{Forgetting the action}

The inclusion of a fiber
\[
X_\Sigma\hookrightarrow EK_N\times_{K_N}X_\Sigma
\]
induces the \emph{nonequivariant limit}
\[
H_{K_N}^*(X_\Sigma;\QQ)
\longrightarrow
H^*(X_\Sigma;\QQ).
\]
On the coefficient ring this is the augmentation
\[
S\longrightarrow\QQ,
\qquad
m\longmapsto0
\quad(m\in M).
\]
For a smooth complete toric variety,
\[
\boxed{
H^*(X_\Sigma;\QQ)
\simeq
H_{K_N}^*(X_\Sigma;\QQ)\otimes_S\QQ
\simeq
\frac{PP^*(\Sigma;\QQ)}{S_+\,PP^*(\Sigma;\QQ)},
}
\]
where \(S_+\) is the ideal of positive-degree global polynomials.

\noindent
\textit{Reference.}
This is Cox--Little--Schenck, Proposition~12.4.7 and
Corollary~12.4.8. Conceptually, the relevant Serre spectral sequence
degenerates because a complete simplicial toric variety has no odd
rational cohomology.

In the Stanley--Reisner presentation, setting every \(m\in M\) to zero
imposes
\[
\theta_m
=
\sum_\rho\langle m,u_\rho\rangle x_\rho
=0.
\]
Thus the ideal \(J_\Sigma\) from Chapter~8 reappears automatically:
\[
\frac{\QQ[x_\rho]}{I_\Sigma}
\quad\longrightarrow\quad
\frac{\QQ[x_\rho]}{I_\Sigma+J_\Sigma}.
\]
In degree two, the same statement says
\[
\boxed{
\text{piecewise linear functions}
\big/
\text{global linear functions}
\;\simeq\;
H^2(X_\Sigma;\QQ),
}
\]
which is the support-function description of divisor classes in yet
another guise.

For \(\PP^1\), the calculation above makes this visible at once:
\[
\frac{\QQ[t,\eta]}{\langle\eta(\eta-t)\rangle}
\otimes_{\QQ[t]}\QQ
\simeq
\frac{\QQ[\eta]}{\langle\eta^2\rangle}
=
H^*(\PP^1;\QQ).
\]

\begin{dictionarytable}
The difference between the two cohomology rings can now be summarized
without any new generators.

\medskip

\begin{center}
\small
\begin{tabularx}{\linewidth}{@{}
  >{\bfseries\raggedright\arraybackslash}p{0.27\linewidth}
  >{\raggedright\arraybackslash}X
  >{\raggedright\arraybackslash}X@{}}
\toprule
& \textbf{Equivariant} & \textbf{Ordinary} \\
\midrule
coefficient ring
& \(S=\operatorname{Sym}_\QQ(M_\QQ)\)
& \(\QQ\)
\\
fan presentation
& \(\QQ[x_\rho]/I_\Sigma\)
& \(\QQ[x_\rho]/(I_\Sigma+J_\Sigma)\)
\\
fan functions
& piecewise polynomials
& quotient by global positive-degree polynomials
\\
degree two
& piecewise linear functions
& piecewise linear modulo global linear
\\
\bottomrule
\end{tabularx}
\end{center}
\end{dictionarytable}

\subsection{A full two-dimensional example: \texorpdfstring{\(\PP^2\)}{P2}}

Let the rays of the fan of \(\PP^2\) be generated by
\[
u_0=(-1,-1),
\qquad
u_1=(1,0),
\qquad
u_2=(0,1),
\]
and let
\[
\sigma_0=\Cone{u_1,u_2},
\qquad
\sigma_1=\Cone{u_2,u_0},
\qquad
\sigma_2=\Cone{u_0,u_1}.
\]
Write \(S=\QQ[x,y]\), where \(x,y\) are the dual coordinates on
\(N_\RR\).

An equivariant class is a triple
\[
(f_0,f_1,f_2)\in S^3.
\]
The three common rays impose three congruences. Since \(x\) vanishes on
\(\rho_2\), \(y\) on \(\rho_1\), and \(y-x\) on \(\rho_0\), we obtain
\[
\boxed{
\begin{aligned}
f_1-f_0&\in(x),\\
f_2-f_0&\in(y),\\
f_2-f_1&\in(y-x).
\end{aligned}}
\]

\begin{figure}[htbp]
\centering
\begin{tikzpicture}[scale=1.05]
  \coordinate (O) at (0,0);
  \coordinate (U1) at (3.0,0);
  \coordinate (U2) at (0,3.0);
  \coordinate (U0) at (-2.25,-2.25);

  \fill[DictionaryBlue!8] (O) -- (U1) -- (U2) -- cycle;
  \fill[SurfaceGreen!10] (O) -- (U2) -- (U0) -- cycle;
  \fill[WarningRed!7] (O) -- (U0) -- (U1) -- cycle;

  \draw[->,thick] (O) -- (3.35,0) node[right] {$u_1$};
  \draw[->,thick] (O) -- (0,3.35) node[above] {$u_2$};
  \draw[->,thick] (O) -- (-2.55,-2.55) node[below left] {$u_0$};

  \node at (1.05,1.05) {$f_0$};
  \node at (-1.1,0.95) {$f_1$};
  \node at (0.85,-1.15) {$f_2$};

  \node[fill=white,inner sep=1.5pt] at (-0.78,-1.48) {\scriptsize \(y-x=0\)};
  \node[fill=white,inner sep=1.5pt] at (1.55,-0.28) {\scriptsize \(y=0\)};
  \node[fill=white,inner sep=1.5pt] at (-0.30,1.55) {\scriptsize \(x=0\)};
\end{tikzpicture}
\caption{Piecewise polynomial data on the fan of \(\PP^2\). Across
each common ray, the difference of the adjacent polynomials is
divisible by the linear form vanishing on that ray.}
\label{fig:p2-piecewise-polynomial}
\end{figure}

Let us solve the compatibility conditions rather than merely state
them. The first two give
\[
f_1=f_0+xg,
\qquad
f_2=f_0+yh
\]
for some \(g,h\in S\). The third condition says
\[
yh-xg\in(y-x).
\]
Setting \(y=x\) shows \(x(h-g)|_{y=x}=0\), hence
\[
h-g\in(y-x).
\]
Write \(h=g+(y-x)k\). Then
\[
\boxed{
(f_0,f_1,f_2)
=
f_0(1,1,1)
+g(0,x,y)
+k(0,0,y(y-x)).
}
\]
Thus \(H_{K_N}^*(\PP^2;\QQ)\) is a free \(S\)-module with one generator
in cohomological degrees \(0,2,4\). This is the equivariant shadow of
the three ordinary cohomology classes
\[
1,\ H,\ H^2
\]
that we used in Section~9.1.

\noindent
\textit{Reference.}
This calculation is Cox--Little--Schenck, Example~12.4.18. We have
spelled out the divisibility step because it is the simplest place to
see fixed-point compatibility become an actual computation.

\begin{warningbox}
Taking the nonequivariant limit is a quotient of the
\emph{piecewise-polynomial ring} by global positive-degree
polynomials. It is not the operation of evaluating each fixed-point
polynomial separately at \(x=y=0\). For example, the class
\((0,x,y)\) above survives as the hyperplane class after the
nonequivariant quotient even though all three displayed polynomials
have zero constant term.
\end{warningbox}

We have now reached the form of equivariant cohomology needed for
localization. A class is a compatible polynomial at every maximal
cone, hence at every torus fixed point. The next question is whether a
\emph{global integral} can also be reconstructed from these local
restrictions.

\begin{guidingquestion}
If an equivariant class is determined by its values at the finitely
many fixed points \(x_\sigma\), can integration over all of
\(X_\Sigma\) also be expressed as a sum of fixed-point contributions?
What must appear in the denominator to compensate for replacing an
\(n\)-dimensional variety by isolated points?
\end{guidingquestion}

\section{Fixed points and Atiyah--Bott localization}
\label{sec:fixed-points-localization}
\index{Atiyah--Bott localization}
\index{fixed point!localization}
\index{tangent weight}
\index{equivariant Euler class}
\index{equivariant integration}

Section~9.2 turned an equivariant cohomology class into a collection
of polynomials
\[
\alpha
\longleftrightarrow
\bigl(\alpha|_{x_\sigma}\bigr)_{\sigma\in\Sigma(n)}
\]
indexed by the torus fixed points.  Localization says something much
stronger: after allowing ourselves to divide by the torus weights,
\emph{integration itself} can be carried out on this finite set.

The denominator has a simple geometric meaning.  A point has lost all
of the tangent directions of \(X\), so its contribution must be divided
by the equivariant Euler class of those missing directions.  For a
toric variety, those tangent directions are read directly from the
maximal cone.

\subsection{The tangent weights at a fixed point}

Let \(\sigma\in\Sigma(n)\) be a maximal cone of a smooth fan and write
\[
\sigma=\Cone{u_{\sigma,1},\ldots,u_{\sigma,n}}.
\]
Smoothness says that the primitive ray generators
\(u_{\sigma,1},\ldots,u_{\sigma,n}\) form a basis of \(N\).  Let
\[
m_{\sigma,1},\ldots,m_{\sigma,n}\in M
\]
be the dual basis.  Equivalently, these are the primitive generators of
the dual cone \(\sigma^\vee\).

The affine chart \(U_\sigma\) is isomorphic to \(\CC^n\), the point
\(x_\sigma\) is its origin, and the characters determined by the
\(m_{\sigma,i}\) describe the \(n\) coordinate directions.  Thus the
fan does not merely tell us where the fixed points are: it tells us the
torus representation on the tangent space at each one.

\begin{proposition}[The equivariant Euler class at a toric fixed point]
\label{prop:toric-fixed-point-euler-class}
Let \(X_\Sigma\) be smooth and complete and let
\(x_\sigma\in X_\Sigma^{K_N}\) correspond to a maximal cone
\(\sigma\).  With the sign convention for equivariant divisor classes
used in Section~9.2,
\[
\boxed{
e_{K_N}(T_{x_\sigma}X_\Sigma)
=
(-1)^n\prod_{i=1}^n m_{\sigma,i}
\in S=\operatorname{Sym}_\QQ(M_\QQ).
}
\]
In particular, this class becomes invertible after we invert the
nonzero linear forms in \(S\).
\end{proposition}

\noindent
\textit{Reference.}
This is the fixed-point calculation behind Cox--Little--Schenck,
Lemma~13.3.5 and Proposition~13.3.6.  The interpretation of the
denominator as an equivariant Euler class is made explicit immediately
after their Theorem~13.3.7 and in Exercise~13.3.2.

\begin{proof}[Conceptual proof sketch]
The invariant divisors corresponding to the rays of \(\sigma\) meet
transversely at \(x_\sigma\).  Their equivariant normal directions
therefore give the \(n\) tangent directions at the fixed point.  The
restriction calculation from Section~9.2 gives
\[
[D_{\rho_i}]_{K_N}|_{x_\sigma}=-m_{\sigma,i}.
\]
The equivariant Euler class is the product of these degree-two classes,
which gives the displayed formula.  This is exactly where the factor
\((-1)^n\) in our localization formula will come from.
\end{proof}

\begin{pauseandthink}
There are two pieces of fan data at \(x_\sigma\): the value
\(\alpha|_{x_\sigma}\), which Section~9.2 records as the polynomial on
\(\sigma\), and the dual basis
\(m_{\sigma,1},\ldots,m_{\sigma,n}\), which records the tangent
directions.  Localization will use exactly these two pieces and no
others.
\end{pauseandthink}

\subsection{Equivariant integration}

Ordinary integration over a smooth complete \(n\)-fold is the
pushforward to a point,
\[
\int_X:H^k(X;\QQ)\longrightarrow H^{k-2n}(\mathrm{pt};\QQ).
\]
There is an equivariant version of the same pushforward,
\[
\int_X^{K_N}:
H_{K_N}^k(X;\QQ)
\longrightarrow
H_{K_N}^{k-2n}(\mathrm{pt};\QQ)=S^{k-2n}.
\]
It remembers the equivariant parameters rather than setting them equal
to zero.  In top degree the answer lies in \(S^0=\QQ\), so for an
equivariant lift of an ordinary top-degree class it agrees with the
ordinary integral after forgetting the action.

To state localization, let
\[
S_{\mathrm{loc}}
:=
S\bigl[\ell^{-1}:0\ne\ell\in M_\QQ\bigr]
\subseteq\Frac(S).
\]
Thus we are allowed to divide by nonzero torus weights.

\begin{theorem}[Atiyah--Bott localization]
\label{thm:atiyah-bott-toric-localization}
Let \(X_\Sigma\) be a smooth complete \(n\)-dimensional toric variety.
For every \(\alpha\in H_{K_N}^*(X_\Sigma;\QQ)\),
\[
\boxed{
\int_{X_\Sigma}^{K_N}\alpha
=
\sum_{\sigma\in\Sigma(n)}
\frac{\alpha|_{x_\sigma}}
{e_{K_N}(T_{x_\sigma}X_\Sigma)}
=
(-1)^n
\sum_{\sigma\in\Sigma(n)}
\frac{\alpha|_{x_\sigma}}
{\prod_{i=1}^n m_{\sigma,i}}.
}
\]
The equality is computed in \(S_{\mathrm{loc}}\), but the sum belongs
to \(S\): all apparent poles cancel.
\end{theorem}

\noindent
\textit{Reference.}
This is Cox--Little--Schenck, Theorem~13.3.7, which is the toric form of
the Atiyah--Bott/Berline--Vergne localization formula.  The localization
of the fixed-point restriction map used in the proof is their
equation~(13.3.3), following the discussion of the localization theorem
on p.~643.

\begin{proof}[Conceptual proof sketch]
Restriction to the fixed points gives
\[
H_{K_N}^*(X_\Sigma)
\longrightarrow
\bigoplus_{\sigma\in\Sigma(n)}S.
\]
Section~9.2 explained why its image consists of compatible polynomial
data.  After the nonzero weights are inverted, the localization theorem
says that this restriction map becomes an isomorphism.  In other
words, the fixed points become a complete basis for equivariant
cohomology once fractions in the weights are allowed.

The class supported at \(x_\sigma\) has restriction
\(e_{K_N}(T_{x_\sigma}X)\) at that fixed point and zero at the others.
Dividing by the Euler class therefore normalizes it to have value one
at \(x_\sigma\).  Decompose \(\alpha\) into these normalized fixed-point
classes and push each one to a point.  The result is exactly the sum in
the theorem.
\end{proof}

At first sight the formula looks perverse.  The left side is a
polynomial, while the individual terms on the right are usually
rational functions.  The cancellation of the denominators is not an
accident: it is the algebraic shadow of the fact that the fixed-point
data came from one global equivariant class.

There is also a useful degree check.  If \(\deg\alpha<2n\), then the
equivariant pushforward would have negative degree and hence must
vanish.  Localization therefore manufactures nontrivial identities
among rational functions of the torus weights.

\subsection{The first localization calculation:
\texorpdfstring{\(\PP^1\)}{P1}}

Return to the notation of Section~9.2.  We have
\[
S=\QQ[t],
\qquad
\sigma_+=\RR_{\ge0},
\qquad
\sigma_-=\RR_{\le0},
\]
and an equivariant class is a pair
\[
(f_+,f_-),
\qquad
f_+-f_-\in(t).
\]
The dual primitive generators are
\[
m_+=t,
\qquad
m_-=-t.
\]
Since \(n=1\), Proposition~\ref{prop:toric-fixed-point-euler-class} gives
\[
e_{K_N}(T_{x_+}\PP^1)=-t,
\qquad
e_{K_N}(T_{x_-}\PP^1)=t.
\]
Therefore localization is simply
\[
\boxed{
\int_{\PP^1}^{K_N}(f_+,f_-)
=
-\frac{f_+}{t}+\frac{f_-}{t}
=
\frac{f_--f_+}{t}.
}
\]
The gluing condition \(f_+-f_-\in(t)\) is precisely what makes the
answer a polynomial.

For the unit class,
\[
\int_{\PP^1}^{K_N}1
=
-\frac1t+\frac1t
=0,
\]
as it must by degree.  More interestingly, recall the equivariant
divisor class
\[
\eta=(0,t),
\qquad
H_{K_N}^*(\PP^1;\QQ)
\simeq
\frac{\QQ[t,\eta]}{\langle\eta(\eta-t)\rangle}.
\]
Then
\[
\boxed{
\int_{\PP^1}^{K_N}\eta
=
0+\frac{t}{t}
=1.
}
\]
After forgetting the equivariant parameter, \(\eta\) becomes the point
class \(H\), and we recover the familiar identity
\[
\int_{\PP^1}H=1.
\]

This example is worth remembering as a miniature of the entire theory:
two local rational expressions, each meaningless at \(t=0\) on its
own, combine to give a perfectly regular global answer.

\subsection{A two-dimensional calculation:
\texorpdfstring{\(\PP^2\)}{P2}}

Now use exactly the fan notation from the end of Section~9.2:
\[
u_0=(-1,-1),
\qquad
u_1=(1,0),
\qquad
u_2=(0,1),
\]
with maximal cones \(\sigma_0,\sigma_1,\sigma_2\) and
\(S=\QQ[x,y]\).  Taking dual bases in the three cones gives the
following fixed-point data.

\begin{center}
\small
\begin{tabularx}{0.94\linewidth}{@{}
  >{\raggedright\arraybackslash}p{0.18\linewidth}
  >{\raggedright\arraybackslash}X
  >{\raggedright\arraybackslash}X@{}}
\toprule
\textbf{Cone}
&
\textbf{Dual generators}
&
\textbf{Euler class}
\\
\midrule
\(\sigma_0=\Cone{u_1,u_2}\)
& \(x,\ y\)
& \(xy\)
\\
\(\sigma_1=\Cone{u_2,u_0}\)
& \(y-x,\ -x\)
& \(x(x-y)\)
\\
\(\sigma_2=\Cone{u_0,u_1}\)
& \(-y,\ x-y\)
& \(y(y-x)\)
\\
\bottomrule
\end{tabularx}
\end{center}

Already the unit class produces an identity:
\[
\begin{aligned}
\int_{\PP^2}^{K_N}1
&=
\frac1{xy}
+\frac1{x(x-y)}
+\frac1{y(y-x)}
\\
&=0.
\end{aligned}
\]
Again this is forced by degree, but localization turns that simple
topological fact into a cancellation among three rational functions.

Now recall the degree-two class from the piecewise-polynomial
calculation of Section~9.2,
\[
h_T=(0,x,y).
\]
Its nonequivariant limit is the hyperplane class \(H\).  Since
\(h_T^2=(0,x^2,y^2)\), localization gives
\[
\begin{aligned}
\int_{\PP^2}^{K_N}h_T^2
&=
\frac{0}{xy}
+\frac{x^2}{x(x-y)}
+\frac{y^2}{y(y-x)}
\\
&=
\frac{x}{x-y}
+\frac{y}{y-x}
=1.
\end{aligned}
\]
Therefore, after taking the nonequivariant limit,
\[
\boxed{
\int_{\PP^2}H^2=1.
}
\]

Notice what we did \emph{not} do: we never multiplied divisor classes
in the ordinary cohomology ring.  The same intersection number was
recovered solely from the three fixed-point restrictions of an
equivariant class and the tangent weights encoded by the fan.

\begin{dictionarytable}
Localization turns the fan description of Section~9.2 into an
integration algorithm.

\medskip

\begin{center}
\small
\begin{tabularx}{\linewidth}{@{}
  >{\bfseries\raggedright\arraybackslash}p{0.25\linewidth}
  >{\raggedright\arraybackslash}X
  >{\raggedright\arraybackslash}X@{}}
\toprule
& \textbf{Global object} & \textbf{Fixed-point data} \\
\midrule
equivariant class
& \(\alpha\in H_{K_N}^*(X_\Sigma)\)
& polynomial \(\alpha|_{x_\sigma}\) on each maximal cone
\\
tangent geometry
& \(TX_\Sigma\)
& \(e_{K_N}(T_{x_\sigma}X)=(-1)^n\prod_i m_{\sigma,i}\)
\\
integration
& \(\displaystyle\int_X^{K_N}\alpha\)
& \(\displaystyle\sum_\sigma
   \frac{\alpha|_{x_\sigma}}{e_{K_N}(T_{x_\sigma}X)}\)
\\
\bottomrule
\end{tabularx}
\end{center}
\end{dictionarytable}

\subsection{Why this is exactly the tool Riemann--Roch needs}

We can already see the shape of the next section.  Let
\[
D=\sum_\rho a_\rho D_\rho
\]
be an invariant Cartier divisor with Cartier data
\(\{m_\sigma\}_{\sigma\in\Sigma(n)}\).  Section~9.2 identified the
restriction of its equivariant divisor class as
\[
[D]_{K_N}|_{x_\sigma}=m_\sigma.
\]
Hence the equivariant Chern character restricts to
\[
\ch_{K_N}\bigl(\mathcal O(D)\bigr)|_{x_\sigma}
=e^{m_\sigma}.
\]

At the same fixed point, the equivariant Chern roots of the tangent
space are, in our convention,
\[
-m_{\sigma,1},\ldots,-m_{\sigma,n}.
\]
Consequently
\[
\frac{
\operatorname{Td}_{K_N}(X)|_{x_\sigma}
}{
e_{K_N}(T_{x_\sigma}X)
}
=
\prod_{i=1}^n
\frac{1}{1-e^{m_{\sigma,i}}}.
\]
Putting the two restrictions into localization gives the formal
fixed-point expression
\[
\boxed{
\int_X^{K_N}
\ch_{K_N}\bigl(\mathcal O(D)\bigr)
\operatorname{Td}_{K_N}(X)
=
\sum_{\sigma\in\Sigma(n)}
\frac{e^{m_\sigma}}
{\prod_{i=1}^n(1-e^{m_{\sigma,i}})}.
}
\]
Here the exponentials are interpreted in the completed equivariant
cohomology ring; no analytic convergence is involved.

\noindent
\textit{Reference.}
This is the local calculation used in Cox--Little--Schenck,
Theorems~13.3.4 and~13.3.7, in their proof of equivariant
Riemann--Roch.  We have isolated it here because it shows exactly why
localization is the correct bridge between the fan and the
Riemann--Roch integral.

The remarkable feature is the denominator.  The Euler-class factor
from localization cancels the numerator in the Todd factor, leaving
\(1-e^{m_{\sigma,i}}\).  These are precisely the geometric-series
denominators one expects when counting lattice points in a cone.  We
will not develop Brion's theorem separately; in the next section we
will use just this observation to reach the chapter's intended
punchline.

\begin{guidingquestion}
The left side of the last formula is an equivariant characteristic
class integral, while the right side resembles a sum of generating
functions for lattice points near the vertices of a polytope.  Why
should its nonequivariant limit be the Euler characteristic of
\(\mathcal O(D)\), and when does that Euler characteristic equal
\(\#(P_D\cap M)\)?
\end{guidingquestion}

\section{Toric Riemann--Roch and lattice points}
\label{sec:toric-rr-lattice-points}
\index{Riemann--Roch!toric}
\index{lattice point!Riemann--Roch count}
\index{equivariant Euler characteristic}
\index{Ehrhart polynomial}

We can now answer the guiding question with almost no new machinery.
There are three quantities that have appeared at different moments in
the notes:
\[
\chi\bigl(X,\mathcal O_X(D)\bigr),
\qquad
\int_X e^{[D]}\operatorname{Td}(X),
\qquad
\#(P_D\cap M).
\]
The first is sheaf cohomology, the second is topology, and the third is
combinatorics.  For a nef toric divisor on a smooth complete toric
variety, all three are the same number.

The equality between the first two is Hirzebruch--Riemann--Roch.  To
identify the first with the third, we need only recall two facts that
we have already encountered in toric geometry: lattice points give a
basis of global sections, and a nef toric divisor has no higher
cohomology.

\subsection{The equivariant Euler characteristic remembers the weights}

Before forgetting the torus action, it is useful to keep track not
only of the dimensions of the cohomology groups but also of their
characters.  If a torus \(T_N\) acts on a line bundle \(L\), define
formally
\[
\chi_{T_N}(X,L)
=
\sum_{i\geq0}(-1)^i\ch_{T_N}H^i(X,L).
\]
Thus the ordinary Euler characteristic remembers only the total
dimension, while the equivariant Euler characteristic remembers the
weight of every torus eigenspace.

For an invariant divisor \(D\), the monomial basis of global sections
is
\[
H^0\bigl(X_\Sigma,\mathcal O(D)\bigr)
=
\bigoplus_{m\in P_D\cap M}\CC\,\character{m}.
\]
If \(D\) is nef and \(X_\Sigma\) is complete, Demazure vanishing gives
\[
H^i\bigl(X_\Sigma,\mathcal O(D)\bigr)=0
\qquad(i>0).
\]
Consequently the equivariant Euler characteristic becomes the finite
character sum
\[
\boxed{
\chi_{T_N}\bigl(X_\Sigma,\mathcal O(D)\bigr)
=
\sum_{m\in P_D\cap M}e^m.
}
\]
Here \(e^m\) is notation for the character of weight \(m\), written in
the exponential notation used by the equivariant Chern character.

\noindent
\textit{References.}
The monomial basis of global sections is Cox--Little--Schenck,
Proposition~4.3.3.  The higher-cohomology statement is their
Theorem~9.2.3 (Demazure vanishing).  Completeness makes
\(|\Sigma|=N_\RR\) convex, so the hypotheses of that theorem are
automatic once \(D\) is nef.

Now compare this finite sum with the expression at the end of
Section~9.3.  Equivariant Riemann--Roch and localization give
\[
\chi_{T_N}\bigl(X_\Sigma,\mathcal O(D)\bigr)
=
\sum_{\sigma\in\Sigma(n)}
\frac{e^{m_\sigma}}
{\prod_{i=1}^n(1-e^{m_{\sigma,i}})}.
\]
Combining the two descriptions gives the striking identity
\begin{equation}
\boxed{
\sum_{m\in P_D\cap M}e^m
=
\sum_{\sigma\in\Sigma(n)}
\frac{e^{m_\sigma}}
{\prod_{i=1}^n(1-e^{m_{\sigma,i}})}.
}
\label{eq:toric-rr-character-identity}
\end{equation}

When \(D\) is ample, each term on the right has a direct polyhedral
meaning.  The Cartier datum \(m_\sigma\) is the vertex of \(P_D\)
corresponding to the maximal cone \(\sigma\), and the vectors
\(m_{\sigma,1},\ldots,m_{\sigma,n}\) generate the tangent cone of
\(P_D\) at that vertex.  Formally expanding geometric series gives
\[
\frac{e^{m_\sigma}}
{\prod_i(1-e^{m_{\sigma,i}})}
=
e^{m_\sigma}
\prod_i\bigl(1+e^{m_{\sigma,i}}+e^{2m_{\sigma,i}}+\cdots\bigr).
\]
Thus a fixed-point contribution is the generating function of lattice
points in the tangent cone at one vertex.  The infinite tails cancel
when the vertex contributions are added, leaving the finite generating
function of the polytope itself.

\begin{pauseandthink}
This is the precise bridge promised at the end of Section~9.3:
\[
\boxed{
\begin{gathered}
\text{torus fixed point}
\longleftrightarrow
\text{vertex of }P_D,
\\[1mm]
\text{tangent weights}
\longleftrightarrow
\text{edge directions at that vertex}.
\end{gathered}
}
\]
We are seeing one of Brion's identities, but we do not need a separate
theory of rational generating functions here.  For us its role is to
explain geometrically why the localization denominators look exactly
like geometric-series denominators.
\end{pauseandthink}

\subsection{The one-dimensional identity in full}

Let us watch the cancellation happen on \(\PP^1\).  Write the rays of
the fan as
\[
\sigma_+=\RR_{\geq0},
\qquad
\sigma_-=\RR_{\leq0},
\]
and take
\[
D=dD_-,\qquad d\geq0,
\]
where \(D_-\) is the invariant divisor belonging to the ray generated
by \(-1\).  The associated polytope is
\[
P_D=[0,d],
\]
and its Cartier data are
\[
m_+=0,
\qquad
m_-=dt.
\]
The primitive generators of the two dual cones are \(t\) and \(-t\),
respectively.  Therefore the right side of
\eqref{eq:toric-rr-character-identity} is
\[
\frac{1}{1-e^t}
+
\frac{e^{dt}}{1-e^{-t}}.
\]
Put \(q=e^t\).  Then
\[
\begin{aligned}
\frac{1}{1-q}
+
\frac{q^d}{1-q^{-1}}
&=
\frac{1}{1-q}
-
\frac{q^{d+1}}{1-q}
\\
&=
\frac{1-q^{d+1}}{1-q}
\\
&=
1+q+\cdots+q^d.
\end{aligned}
\]
In other words,
\[
\boxed{
\frac{1}{1-e^t}
+
\frac{e^{dt}}{1-e^{-t}}
=
1+e^t+\cdots+e^{dt}.
}
\]

The left side is a sum of two infinite-cone contributions, one from
each fixed point.  The right side contains exactly the \(d+1\) lattice
points of the interval \([0,d]\).  Finally forgetting the torus action
means sending every character to its dimension, or informally taking
\(e^t\to1\).  Although the two fixed-point fractions individually have
poles there, their sum has the perfectly regular limit
\[
\chi\bigl(\PP^1,\mathcal O(d)\bigr)=d+1.
\]

\noindent
\textit{Reference.}
This is the one-dimensional instance of the calculation in
Cox--Little--Schenck, Example~13.3.9.  Their Theorem~13.3.4 gives the
general fixed-point character decomposition.

\subsection{The toric Riemann--Roch counting formula}

We can now forget the equivariant parameters and state the numerical
payoff of the chapter.

\begin{theorem}[Toric Riemann--Roch as lattice-point counting]
\label{thm:toric-rr-lattice-count}
Let \(X_\Sigma\) be a smooth complete toric variety and let \(D\) be a
nef invariant Cartier divisor.  Then
\[
\boxed{
\#(P_D\cap M)
=
\chi\bigl(X_\Sigma,\mathcal O(D)\bigr)
=
\int_{X_\Sigma}
e^{[D]}
\prod_{\rho\in\Sigma(1)}
\frac{[D_\rho]}{1-e^{-[D_\rho]}}.
}
\]
In particular, the formula holds when \(D\) is ample.
\end{theorem}

\noindent
\textit{References.}
This is the direct combination of Cox--Little--Schenck,
Proposition~4.3.3, Theorem~9.2.3, Theorem~13.1.6, and
Theorem~13.3.8.  The last result is their toric
Hirzebruch--Riemann--Roch theorem.

\begin{proof}[Conceptual proof]
Hirzebruch--Riemann--Roch gives
\[
\chi\bigl(X_\Sigma,\mathcal O(D)\bigr)
=
\int_{X_\Sigma}e^{[D]}\operatorname{Td}(X_\Sigma).
\]
The toric Euler sequence turns the Todd class into the displayed
product over invariant divisors.  Since \(D\) is nef, Demazure
vanishing kills every higher cohomology group, so
\[
\chi\bigl(X_\Sigma,\mathcal O(D)\bigr)
=h^0\bigl(X_\Sigma,\mathcal O(D)\bigr).
\]
Finally, the characters \(\character{m}\), for \(m\in P_D\cap M\),
form a basis of global sections.  Hence
\[
h^0\bigl(X_\Sigma,\mathcal O(D)\bigr)=\#(P_D\cap M),
\]
as claimed.
\end{proof}

There is a useful distinction in the theorem.  Riemann--Roch itself
always computes the Euler characteristic of a line bundle on our
smooth complete toric variety.  The positivity hypothesis is used
only in the second step, when we identify that alternating sum with
\(h^0\) and hence with an honest lattice-point count.

\begin{warningbox}
Do not silently replace \(\chi(X,\mathcal O(D))\) by \(h^0(X,\mathcal
O(D))\) for an arbitrary divisor.  For negative divisors the higher
cohomology can be nonzero.  The example
\(\chi(\PP^1,\mathcal O(-2))=-1\) from Section~9.1 is the simplest
warning: an Euler characteristic can be negative, while a number of
lattice points cannot.
\end{warningbox}

\subsection{A corner cut: two computations on
\texorpdfstring{\(\Bl_p(\PP^2)\)}{Blp(P2)}}

Let
\[
X=\Bl_p(\PP^2)
\]
at a torus fixed point, with the notation from Section~8.4.  Thus
\[
H^*(X;\ZZ)
\simeq
\frac{\ZZ[H,E]}{\langle HE,H^2+E^2\rangle},
\qquad
\int_XH^2=1,
\]
where \(H\) is the pullback of a line and \(E\) is the exceptional
curve.  The canonical class is
\[
K_X=-3H+E.
\]

Take
\[
D=dH-mE,
\qquad
d\geq m\geq0.
\]
This divisor is nef.  In the ray notation from Section~8.4 it can be
represented by
\[
D=dD_0-mD_E,
\]
so its polytope is cut out by
\[
x\geq0,
\qquad
y\geq0,
\qquad
m\leq x+y\leq d.
\]
It is the degree-\(d\) triangle with the corner
\(x+y<m\) removed.

\begin{figure}[htbp]
\centering
\begin{tikzpicture}[scale=0.86]
  \fill[DictionaryBlue!10]
    (1.6,0)--(5.0,0)--(0,5.0)--(0,1.6)--cycle;
  \draw[thick,DictionaryBlue]
    (1.6,0)--(5.0,0)--(0,5.0)--(0,1.6)--cycle;
  \draw[dashed,PausePurple,thick] (0,0)--(1.6,0)--(0,1.6)--cycle;
  \draw[->] (-0.25,0)--(5.65,0);
  \draw[->] (0,-0.25)--(0,5.65);
  \node[below] at (1.6,0) {$(m,0)$};
  \node[below] at (5.0,0) {$(d,0)$};
  \node[left] at (0,1.6) {$(0,m)$};
  \node[left] at (0,5.0) {$(0,d)$};
  \node[DictionaryBlue] at (2.65,2.25) {$P_D$};
  \node[PausePurple,align=center] at (0.95,0.72)
    {corner\\removed};
\end{tikzpicture}
\caption{The polytope of \(D=dH-mE\) is obtained from the degree-\(d\)
triangle by cutting off the lattice corner of size \(m\).}
\label{fig:blowup-rr-corner-cut}
\end{figure}

First compute by Riemann--Roch.  From the intersection relations,
\[
D^2=d^2-m^2,
\]
and
\[
\begin{aligned}
D\cdot(D-K_X)
&=(dH-mE)\bigl((d+3)H-(m+1)E\bigr)
\\
&=d(d+3)-m(m+1).
\end{aligned}
\]
The surface formula from Section~9.1 therefore gives
\[
\begin{aligned}
\chi\bigl(X,\mathcal O_X(dH-mE)\bigr)
&=1+\frac12\bigl[d(d+3)-m(m+1)\bigr]
\\
&=
\binom{d+2}{2}-\binom{m+1}{2}.
\end{aligned}
\]

Now compute from the polygon.  The full degree-\(d\) triangle has
\[
\binom{d+2}{2}
\]
lattice points.  The discarded lattice points satisfy
\(x,y\geq0\) and \(x+y<m\).  There are
\[
1+2+\cdots+m
=
\binom{m+1}{2}
\]
of them.  Hence
\[
\boxed{
\#(P_D\cap M)
=
\binom{d+2}{2}-\binom{m+1}{2}
=
\chi\bigl(X,\mathcal O_X(dH-mE)\bigr).
}
\]

The calculation has a familiar algebraic interpretation as well:
sections of \(\mathcal O_{\PP^2}(d)\) are degree-\(d\) monomials, and
passing to \(dH-mE\) imposes vanishing of order at least \(m\) at the
blown-up point.  The monomials removed by that vanishing condition are
exactly the lattice points in the cut-off corner.

\begin{threedictionaries}
The corner-cut calculation is the chapter in miniature.

\medskip

\begin{center}
\small
\begin{tabularx}{\linewidth}{@{}
  >{\raggedright\arraybackslash}X
  >{\raggedright\arraybackslash}X
  >{\raggedright\arraybackslash}X@{}}
\toprule
\textbf{Algebraic geometry}
&
\textbf{Topology}
&
\textbf{Polytope}
\\
\midrule
\(\mathcal O_X(dH-mE)\)
& \(dH-mE\in H^2(X)\)
& \(m\leq x+y\leq d\)
\\
vanish to order \(m\)
& exceptional class \(E^2=-1\)
& cut off a corner of size \(m\)
\\
Euler characteristic
& HRR integral
& number of lattice points
\\
\bottomrule
\end{tabularx}
\end{center}
\end{threedictionaries}

\subsection{Dilating the polytope}

There is one final consequence worth recording.  If \(D\) is ample
and \(P=P_D\), then \(kD\) corresponds to the dilated polytope
\(kP\).  Applying Theorem~\ref{thm:toric-rr-lattice-count} to every
integer \(k\geq0\) gives
\[
\boxed{
\#(kP\cap M)
=
\int_X e^{k[D]}\operatorname{Td}(X).
}
\]
Because cohomology vanishes above degree \(2n\), the exponential on
the right truncates.  The right side is therefore a polynomial in
\(k\) of degree at most \(n\).  Thus, for the smooth lattice polytopes
arising here, Riemann--Roch gives a geometric explanation for the
polynomial behavior of lattice-point counts under dilation.

We will not develop the general theory of Ehrhart polynomials or the
volume-polynomial and Khovanskii--Pukhlikov refinements here.  The
important lesson for these notes is already visible: characteristic
classes convert the geometry of a toric variety into arithmetic data
of its polytope.

\begin{tcolorbox}[
  colback=PredictionTeal!5,
  colframe=PredictionTeal,
  title={The punchline},
  fonttitle=\bfseries
]
For a smooth projective toric variety, the same integer can now be
computed in three genuinely different ways:
\[
\boxed{
\begin{aligned}
\text{sheaf theory:}\quad
&\chi\bigl(X,\mathcal O(D)\bigr),
\\[1mm]
\text{cohomology:}\quad
&\int_X e^{[D]}\operatorname{Td}(X),
\\[1mm]
\text{combinatorics:}\quad
&\#(P_D\cap M).
\end{aligned}}
\]
Equivariant cohomology and localization explain why the bridge between
the last two is controlled one fixed point---equivalently, one vertex
of the polytope---at a time.
\end{tcolorbox}

\clearpage
\section*{What this chapter has accomplished}
\addcontentsline{toc}{section}{What this chapter has accomplished}
\markboth{What this chapter has accomplished}{What this chapter has accomplished}

Chapter~8 gave us the cohomology ring of a smooth complete toric
variety.  The first lesson of this chapter was that this ring can be
used to compute an algebro-geometric invariant.  For a line bundle
\(L\) on a smooth complete variety, Hirzebruch--Riemann--Roch says
\[
\boxed{
\chi(X,L)
=
\int_X\ch(L)\operatorname{Td}(X).
}
\]
For \(L=\mathcal O_X(D)\), the Chern character is \(e^{[D]}\), while
the generalized toric Euler sequence turns the Todd class into the
explicit product
\[
\boxed{
\operatorname{Td}(X_\Sigma)
=
\prod_{\rho\in\Sigma(1)}
\frac{[D_\rho]}{1-e^{-[D_\rho]}}.
}
\]
On a surface this recovers the particularly concrete formula
\[
\chi\bigl(X,\mathcal O_X(D)\bigr)
=
1+\frac12D\cdot(D-K_X).
\]

Equivariant cohomology then put the torus action back into the story.
For the compact torus \(K_N\subset T_N\), its coefficient ring is
\[
H_{K_N}^*(\mathrm{pt};\QQ)
\simeq
\operatorname{Sym}_\QQ(M_\QQ)=:S,
\]
and a smooth complete toric variety is equivariantly formal.  Most
importantly, its equivariant cohomology has a fan-theoretic model:
\[
\boxed{
H_{K_N}^*(X_\Sigma;\QQ)
\simeq
PP^*(\Sigma;\QQ).
}
\]
Thus an equivariant class is a polynomial on each cone, with the
polynomials agreeing on common faces.  In degree two these are the
piecewise linear functions already familiar from toric divisors.
Forgetting the action amounts to setting the positive-degree elements
of \(S\) equal to zero; this recovers ordinary cohomology.

Localization showed why the piecewise-polynomial description is so
powerful.  A maximal cone \(\sigma\) determines both a fixed point
\(x_\sigma\) and the tangent weights
\(m_{\sigma,1},\ldots,m_{\sigma,n}\).  After those weights are made
invertible, integration becomes a finite sum:
\[
\boxed{
\int_{X_\Sigma}^{K_N}\alpha
=
(-1)^n
\sum_{\sigma\in\Sigma(n)}
\frac{\alpha|_{x_\sigma}}
     {\prod_i m_{\sigma,i}}.
}
\]
The individual summands are rational functions, but the apparent
poles cancel because the local data come from one global equivariant
class.

Finally, Riemann--Roch, positivity, and the monomial basis of toric
sections brought the entire chapter back to the polytope.  If \(D\)
is nef and Cartier, then Demazure vanishing gives
\[
\boxed{
\chi\bigl(X_\Sigma,\mathcal O(D)\bigr)
=
\int_{X_\Sigma}e^{[D]}\operatorname{Td}(X_\Sigma)
=
\#(P_D\cap M).
}
\]
Before forgetting the action, the same statement remembers each
lattice point as a character, and localization decomposes the answer
into one contribution from each vertex of the polytope.  The fan,
equivariant topology, and lattice-point combinatorics are therefore
not three parallel descriptions: they are three stages of one
calculation.

The challenge exercise reverses an argument we used earlier in the
notes.  Chapter~7 invoked Pick's theorem to extract geometry from
lattice polygons.  We now have enough machinery to derive Pick's
theorem itself from toric Riemann--Roch.

\begin{exercises}[Comprehension exercises]
\markboth{Comprehension exercises}{Comprehension exercises}

\item Let \(A\) and \(B\) be the two ruling classes on
\(X=\PP^1\times\PP^1\), so that
\[
A^2=B^2=0,
\qquad
A\cdot B=1,
\qquad
K_X=-2A-2B.
\]
For \(D=aA+bB\), use surface Riemann--Roch to prove
\[
\chi\bigl(X,\mathcal O_X(a,b)\bigr)=(a+1)(b+1).
\]
When \(a,b\geq0\), recover the same number by counting lattice points
in the rectangle \([0,a]\times[0,b]\).

\item Let \(S,F\in H^2(\mathbb F_r;\ZZ)\) be the section and fiber
classes from the Chapter~8 challenge, with
\[
S^2=-r,
\qquad
F^2=0,
\qquad
S\cdot F=1.
\]
Use the four invariant boundary divisors to show that
\[
K_{\mathbb F_r}=-2S-(r+2)F.
\]
Then compute \(\chi(\mathbb F_r,\mathcal O(aS+bF))\) directly from
surface Riemann--Roch.  Which part of your answer remembers the
integer \(r\), even though Chapter~8 showed that the abstract
cohomology ring remembers only its parity?

\item Let \(S^1\) act on \(S^3\subset\CC^2\) by scalar
multiplication.  The action is free.  Use the Borel construction to
compute
\[
H_{S^1}^*(S^3;\QQ).
\]
Now give the same underlying space the trivial \(S^1\)-action and
compute its equivariant cohomology again.  Explain precisely where
the two answers record the difference between the actions.

\item Let \(T=(\CC^*)^2\) act in the standard way on
\(X=\PP^1\times\PP^1\).  Without finding a presentation of the whole
ring, use equivariant formality to determine the rank of
\(H_T^*(X;\QQ)\) as a module over
\[
H_T^*(\mathrm{pt};\QQ)=\QQ[t_1,t_2].
\]
Then explain algebraically how setting \(t_1=t_2=0\) recovers a
four-dimensional ordinary cohomology ring.  Why should this rank also
equal the number of torus fixed points?

\item Consider the fan of \(\PP^1\times\PP^1\) with rays generated by
\[
(1,0),\quad(0,1),\quad(-1,0),\quad(0,-1).
\]
Describe an element of \(PP^*(\Sigma;\QQ)\) as four polynomials, one
on each quadrant, and write the compatibility condition across each
ray.  Specialize your description to degree two and explain how the
resulting piecewise linear functions recover torus-invariant Cartier
divisor classes.  What changes after quotienting by global linear
functions?

\item On each factor of \(\PP^1\times\PP^1\), use the equivariant
class \(\eta=(0,t)\) from Section~9.3.  Let \(\eta_1,\eta_2\) be its
pullbacks from the two factors, with equivariant parameters
\(t_1,t_2\).  Use the four fixed points and localization to prove
\[
\int_{\PP^1\times\PP^1}^{T}
\eta_1\eta_2=1.
\]
Then use localization again and watch the two contributions in the
other factor cancel to show
\[
\int_X^T\eta_1^2
=
\int_X^T\eta_2^2
=0.
\]
Recover the ordinary intersection number
\((A+B)^2=2\).

\item For \(d\geq0\), apply Riemann--Roch on \(\PP^1\) to
\(\mathcal O_{\PP^1}(-d-2)\).  Show that
\[
\chi\bigl(\PP^1,\mathcal O(-d-2)\bigr)=-d-1.
\]
Since this line bundle has no nonzero global sections, determine
\(h^1(\PP^1,\mathcal O(-d-2))\).  Explain why this example shows that
the nef hypothesis in the lattice-point version of toric
Riemann--Roch cannot simply be discarded.

\item Let a nef divisor on \(\mathbb F_r\) have polygon
\[
P=\Conv{(0,0),(a,0),(a+rb,b),(0,b)},
\qquad a,b,r\geq0.
\]
Count its lattice points directly by slicing at integer heights
\(y=0,1,\ldots,b\), and prove
\[
\#(P\cap M)
=
(a+1)(b+1)+\frac r2b(b+1).
\]
Now compute the area and the four lattice edge lengths of \(P\).
Use the area/intersection and edge-length/intersection dictionaries
from Chapter~5 together with surface Riemann--Roch to recover exactly
the same answer.

\end{exercises}

\clearpage
\section*{Challenge exercise}
\addcontentsline{toc}{section}{Challenge exercise}
\markboth{Challenge exercise}{Challenge exercise}

\noindent
\textbf{Riemann--Roch proves Pick's theorem.}

\medskip

Let \(P\subset M_\RR\simeq\RR^2\) be a two-dimensional lattice
polygon.  Write
\[
A(P)=\operatorname{Area}(P),
\qquad
B(P)=|\partial P\cap M|,
\qquad
I(P)=|\operatorname{Int}(P)\cap M|.
\]
Pick's theorem is the identity
\[
\boxed{
A(P)=I(P)+\frac12B(P)-1.
}
\]
We used this formula in Chapter~7.  The goal of this challenge is to
reverse the logic: derive it from the toric Riemann--Roch theorem of
this chapter.  Begin with a smooth polygon, where every piece of the
dictionary is already available, and then use a toric resolution to
remove the smoothness assumption.

\begin{enumerate}[label=\textbf{Hint \arabic*.},wide=0pt]

\item \textbf{Turn the polygon into a line bundle.}

First assume that \(P\) is smooth, let \(X=X_P\) be its smooth
projective toric surface, and choose the ample torus-invariant divisor
\(D\) with
\[
P_D=P.
\]
Explain why
\[
P_{kD}=kP
\qquad(k\geq0)
\]
and why Theorem~\ref{thm:toric-rr-lattice-count} gives
\[
L_P(k):=\#(kP\cap M)
=
\chi\bigl(X,\mathcal O_X(kD)\bigr).
\]

\item \textbf{Write the surface Riemann--Roch polynomial.}

Apply the surface formula from Section~9.1 to \(kD\).  Prove
\[
\boxed{
L_P(k)
=
1+\frac12k^2D^2-\frac12k\,D\cdot K_X.
}
\]
Why is the constant term equal to \(1\)?  Give both interpretations:
\(L_P(0)=1\) on the polytope side, and
\(\chi(X,\mathcal O_X)=1\) on the toric side.

\item \textbf{Recognize the quadratic coefficient as area.}

Use the area/self-intersection theorem from Chapter~5 to replace the
self-intersection by
\[
D^2=2A(P).
\]
The preceding polynomial becomes
\[
L_P(k)
=
A(P)k^2-\frac12(D\cdot K_X)k+1.
\]
Notice that no interior lattice-point count has entered the argument.

\item \textbf{Recognize the linear coefficient as boundary length.}

Write the toric boundary as
\[
-K_X=\sum_{\rho\in\Sigma(1)}D_\rho.
\]
If \(F_\rho\) is the edge of \(P\) corresponding to \(\rho\), use
the edge-length/intersection dictionary from Chapter~5 to show
\[
-D\cdot K_X
=
\sum_\rho D\cdot D_\rho
=
\sum_\rho\length_M(F_\rho).
\]
Now prove the elementary boundary identity
\[
\boxed{
\sum_{\substack{F\preceq P\\ \dim F=1}}
\length_M(F)=B(P).
}
\]
Be careful about the vertices: summing \(|F\cap M|\) would count every
vertex twice, whereas summing \(|F\cap M|-1\) counts the cyclic
boundary exactly once.

Conclude that
\[
\boxed{
L_P(k)
=
A(P)k^2+\frac12B(P)k+1.
}
\]
This is the Ehrhart polynomial of a smooth lattice polygon, obtained
here from Riemann--Roch rather than from Ehrhart theory.

\item \textbf{Set \(k=1\).}

Since every lattice point of \(P\) is either interior or on the
boundary,
\[
L_P(1)=I(P)+B(P).
\]
Substitute \(k=1\) into the Riemann--Roch polynomial and show
\[
I(P)+B(P)
=
A(P)+\frac12B(P)+1.
\]
Rearrange.  You have proved
\[
\boxed{
A(P)=I(P)+\frac12B(P)-1,
}
\]
Pick's theorem for smooth lattice polygons.

\item \textbf{Recover the Chapter~7 anticanonical calculation.}

Suppose \(P=P_{-K_X}\) is a smooth reflexive polygon.  Recall that
the origin is its unique interior lattice point, so \(I(P)=1\).
Use the formula you have just proved to recover
\[
2A(P)=B(P).
\]
Then use \((-K_X)^2=2A(P)\) to obtain
\[
\boxed{
(-K_X)^2=B(P),
}
\]
the degree formula used in Section~7.5.  Explain why the logical
direction is now the reverse of the argument in Chapter~7.

\item \textbf{Remove the smoothness assumption.}

Now let \(P\) be an arbitrary two-dimensional lattice polygon.  Its
normal fan \(\Sigma_P\) need not be smooth.  Choose a smooth complete
refinement \(\widetilde\Sigma\) and let
\[
\pi:X_{\widetilde\Sigma}\longrightarrow X_{\Sigma_P}
\]
be the corresponding toric resolution.  If \(A_P\) is the ample
Cartier divisor determined by \(P\) on \(X_{\Sigma_P}\), set
\[
D=\pi^*A_P.
\]
Use the pullback description of basepoint-free toric divisors from
Chapter~5 to explain why \(D\) is nef and
\[
P_D=P.
\]
Thus Theorem~\ref{thm:toric-rr-lattice-count} still applies on the
\emph{smooth} surface \(X_{\widetilde\Sigma}\).

The new rays of \(\widetilde\Sigma\) correspond to curves contracted
by \(\pi\), so \(D\) has intersection zero with those curves.  The
uncontracted boundary curves correspond to the original edges of
\(P\), and their intersections with \(D\) are the lattice lengths of
those edges.  Deduce again that
\[
D^2=2A(P),
\qquad
-D\cdot K_{X_{\widetilde\Sigma}}=B(P).
\]
Repeat Hints 2--5 on the resolution and conclude that Pick's theorem
holds for every lattice polygon.

\item \textbf{One theorem, three languages.}

Summarize the proof by completing the following dictionary in your
own words:
\[
\boxed{
\begin{gathered}
D^2
\longleftrightarrow
2A(P),
\qquad
-D\cdot K_X
\longleftrightarrow
B(P),
\qquad
\chi(\mathcal O_X(D))
\longleftrightarrow
I(P)+B(P).
\end{gathered}}
\]
Which input is topological, which is intersection-theoretic, and
which is combinatorial?  Where, exactly, does positivity enter the
proof?  Where, exactly, is smoothness used, and how did the resolution
remove it?

\end{enumerate}

\begin{threedictionaries}
The proof of Pick's theorem is the chapter's three-way dictionary in
its most compact form.

\medskip

\begin{center}
\small
\begin{tabularx}{\linewidth}{@{}
  >{\raggedright\arraybackslash}X
  >{\raggedright\arraybackslash}X
  >{\raggedright\arraybackslash}X@{}}
\toprule
\textbf{Intersection theory}
& \textbf{Riemann--Roch}
& \textbf{Lattice polygon}
\\
\midrule
\(D^2\)
& quadratic coefficient
& \(2\operatorname{Area}(P)\)
\\
\(-D\cdot K_X\)
& linear coefficient
& \(|\partial P\cap M|\)
\\
\(\chi(\mathcal O_X)\)
& constant coefficient
& \(1\)
\\
\(\chi(\mathcal O_X(D))\)
& value at \(k=1\)
& \(|P\cap M|\)
\\
\bottomrule
\end{tabularx}
\end{center}
\end{threedictionaries}

The proof explains why Pick's formula has exactly three terms.  On a
surface, Riemann--Roch has a quadratic intersection term, a linear
canonical term, and the constant \(\chi(\mathcal O_X)=1\); the toric
dictionary turns them into area, boundary length, and \(1\).

\clearpage
\section*{Notes and references}
\addcontentsline{toc}{section}{Notes and references}
\markboth{Notes and references}{Notes and references}

The Chern-class calculation for a smooth toric variety follows from
the generalized Euler sequence; see Cox--Little--Schenck,
Proposition~13.1.2.  Their Theorem~13.1.6 gives the product formula for
the Todd class, and Example~13.1.7 gives the surface Riemann--Roch
calculation used repeatedly in this chapter.  Hirzebruch--Riemann--
Roch for line bundles on smooth complete toric varieties is stated in
their Theorem~13.3.8.

Our equivariant cohomology discussion follows Cox--Little--Schenck,
\S12.4.  Definition~12.4.6 introduces the Borel construction,
Proposition~12.4.7 and Corollary~12.4.8 describe the passage back to
ordinary cohomology, and Theorem~12.4.14 gives the equivariant
Stanley--Reisner presentation.  The piecewise-polynomial viewpoint is
developed in the discussion surrounding their exact
sequence~(12.4.18).  In degree two it recovers the support functions
and Cartier data already used throughout these notes.

The fixed-point formula in Section~9.3 is the toric form of the
Atiyah--Bott/Berline--Vergne localization theorem.  In
Cox--Little--Schenck, Lemma~13.3.5 and Proposition~13.3.6 compute the
toric fixed-point weights, while Theorem~13.3.7 gives the localization
formula itself.  Their Theorem~13.3.4 and Example~13.3.9 explain the
equivariant Riemann--Roch character calculation that underlies our
one-dimensional example.

The final lattice-point formula combines three earlier results:
global sections are indexed by \(P_D\cap M\) by
Cox--Little--Schenck, Proposition~4.3.3; higher cohomology vanishes for
nef toric divisors by their Theorem~9.2.3; and toric
Hirzebruch--Riemann--Roch is Theorem~13.3.8.  The area and edge-length
intersection formulas used in the challenge were proved earlier in
Chapter~5 and correspond to Cox--Little--Schenck,
Proposition~10.5.6 and Proposition~6.3.8, respectively.

We have intentionally not developed the volume-polynomial,
Khovanskii--Pukhlikov, or general Ehrhart-theoretic refinements of
Chapter~13 of Cox--Little--Schenck.  The point here was narrower: to
see equivariant cohomology and localization turn Riemann--Roch into a
concrete toric computation, and then to watch that computation become
lattice-point arithmetic.  The next chapter changes direction from
invariants of a fixed toric variety to birational geometry: varying a
Cox linearization will move us through the secondary fan and make
contractions, flips, and the toric minimal model program visible.
